\subsection{Respuesta de los escenarios a la matriz y estado de los objetivos}
\label{sec:matriz-escenarios}
Las corridas de la revisión 01 cuantifican E01: C1 reduce la demanda total de enfriamiento un 1,95\,\%, C2 un 0,51\,\% y C3 un 2,37\,\% respecto al CB. Este indicador comparativo no se transforma en cumplimiento de E02, que requiere energía operacional y referencia EDGE. Tampoco permite resolver C01 sin evaluar condiciones térmicas ocupadas.

C1 y C3 incorporan medidas relacionadas con E04 y P01, pero falta verificar su protección solar eficaz, área neta de abertura e incidencia lumínica anual. C2 y C3 modifican la absortancia de cubierta de 0,40 a 0,20; M01 permanece no evaluable porque no hay cantidades y factores ambientales del producto que acrediten menor carbono incorporado. Un cambio óptico no sustituye una selección de materiales sostenibles.

C3 revisión 02 responde a las brechas de envolvente e iluminación mediante nuevos parámetros y una corrida independiente (sección~\ref{sec:c3-revision02}). Su reducción térmica frente al CB es del 22,41\,\%. El cumplimiento completo de E03 y E05 exige, además, detalles constructivos, fotometría y solución de control para las seis zonas sin ventanas exteriores. Los indicadores de agua, materiales, equipos y operación se vinculan con el plan de cierre; una especificación propuesta no cambia automáticamente su estado a cumplimiento.

\begin{table}[H]\centering
\caption{Estado de desarrollo de los objetivos específicos}
{\small\singlespacing\renewcommand{\arraystretch}{1.3}\setlength{\tabcolsep}{4pt}
\begin{tabular}{>{\raggedright\arraybackslash}p{1.6cm}>{\raggedright\arraybackslash}p{3.2cm}>{\raggedright\arraybackslash}p{9.1cm}}\toprule
Objetivo & Estado documental & Evidencia y condición pendiente \\ \midrule
OE1 & Matriz configurada & 16 indicadores, cinco marcos, reglas y fuentes explícitas. Pendiente contraste externo de pertinencia y confirmación de condiciones de aplicación. \\
OE2 & Evaluación parcial & Matriz aplicada al CB, una brecha de iluminación identificada y límites registrados. Faltan datos para completar confort, aire, agua, materiales y otros criterios. \\
OE3 & Propuesta preliminar & Paquete pasivo y acabado reflectante simulados; revisión 02 amplía envolvente e iluminación. Falta selección ambiental demostrada y cierre de los criterios restantes. \\
OE4 & Validación parcial & Demanda térmica comparada en cinco corridas y temperatura operativa ocupada de seis oficinas contrastada entre CB y C3 r02. Se cuantificaron horas con PMV $\pm0,5$; resta separar regímenes de operación y ampliar la cobertura espacial. \\
\bottomrule\end{tabular}}
\fuente{Elaboración propia a partir de las simulaciones y de la matriz de indicadores. Los estados no son calificaciones de certificación.}
\end{table}
